\chapter{Статусная заморозка дефектно-транспортной ветви}

Вторая часть Тома V заменила вопрос о статическом выборе состояния
вопросом о переносе дефекта. Эта смена постановки дала новые точные
результаты: единый дираковский предел, трёхциклические спектральные ветви
$0,\pm1/3$ и одну инвариантную нулевую ветвь на левом канале. Вместе с тем
прямой переход к локализованному ядру потребовал поля спаривания, которого
нет в текущем родителе.

Настоящая глава не добавляет новое поле. Она сопоставляет вторую часть со
всеми прежними майорановскими и колчанными ветвями и фиксирует точный тип
оставшегося разрыва.

\section{Сохранённое положительное ядро}

При заморозке сохраняются следующие результаты:
\begin{enumerate}
 \item наблюдаемый пакет
 \begin{equation}
 H_{15}=Q_L\oplus L_L\oplus u_R\oplus d_R\oplus e_R;
 \end{equation}
 \item переходный носитель $M_{20\times15}(\mathbb C)$, связывающая
 алгебра $M_{35}(\mathbb C)$ и веса углов $4/7,3/7$;
 \item локальный унитарный перенос с дираковским непрерывным пределом;
 \item голономные ветви $0,\pm1/3$;
 \item одна комплексная нулевая ветвь после левого хирального ограничения;
 \item условная нечётная вещественная чётность после майорановского
 ограничения.
\end{enumerate}
Эти утверждения являются кинематическими и спектральными. Они не требуют
объявлять нулевую моду локализованной частицей.

\section{Скрытое смешение двух архитектур}

Старая вихревая ветвь строилась вокруг перенормируемой связи
\begin{equation}
 \Phi N^cN^c,
 \label{eq:v5-freeze-old-majorana-vertex}
\end{equation}
где $N^c$ является правым стерильным состоянием, а $\Phi$ имеет заряд
минус два. Минимальное расширение $U(1)_{B-L}$ действительно даёт
правильные заряды, корневую фазу и сокращение непрерывных аномалий.

Но $N^c$ отсутствует в $H_{15}$. Его добавление означает переход к
\begin{equation}
 H_{16}=H_{15}\oplus N^c,
\end{equation}
и меняет всю родительскую архитектуру:
\begin{equation}
 M_{20\times15}(\mathbb C)\longrightarrow
 M_{20\times16}(\mathbb C),
 \qquad
 M_{35}(\mathbb C)\longrightarrow M_{36}(\mathbb C),
 \label{eq:v5-freeze-h15-h16-carrier}
\end{equation}
\begin{equation}
 \left(\frac47,\frac37\right)
 \longrightarrow
 \left(\frac59,\frac49\right).
 \label{eq:v5-freeze-h15-h16-weights}
\end{equation}
Поэтому условный вихрь \eqref{eq:v5-freeze-old-majorana-vertex} нельзя
вставить в текущий $M_{35}$ без скрытого шестнадцатого состояния.

\begin{nogocriterion}[Запрет смешения $H_{15}$ и $B-L$-вихря]
Нельзя одновременно сохранять носитель $M_{20\times15}(\mathbb C)$, веса
$4/7,3/7$ и использовать вершину $\Phi N^cN^c$. Такая вершина принадлежит
архитектуре $H_{16}$ и требует повторного вывода следа, меры и всех
относительных нормировок.
\end{nogocriterion}

\section{Обратный аудит источников поля спаривания}

Полное меню проекта даёт пять различных частичных результатов.

\subsection{Существующие фазы корня}

Минимальные $SU(5)$, спиновая и намбу-структуры не назначают обязательный
корневой заряд стерильному каналу. Сингулярный квадратный корень
торсионной линии существует на дополнении ядра, но является новым
топологическим сектором или оператором беспорядка.

\subsection{Расширение \texorpdfstring{$B-L$}{B-L}}

Расширение $U(1)_{B-L}$ правильно типизирует $N^c$, $\Phi$ и
\eqref{eq:v5-freeze-old-majorana-vertex}. Однако оно вводит новый
калибровочный множитель, масштаб нарушения, кинетическое смешивание и
матрицу майорановских связей. Кроме того, однородный конденсат несовместим
с корневой голономией; неоднородная конфигурация существует лишь условно
выше невыведенного порога.

\subsection{Семейный колчан}

Трёхузловой колчан точно фиксирует $Y=\Phi I_3$ и алгебраический квадрат
отображения момента. Его KO6-представление существует, но обычный
спектральный след даёт $(\rho^2+r^2)^2$ вместо
$(\rho^2-r^2)^2$. Следовательно, правильный потенциал не следует из
стандартного действия этого конечного оператора.

\subsection{Относительная кривизна и вспомогательное поле}

Отображающий конус выбирает не тот блок, строгая проекция на
$\mathfrak{so}(3)$ равна нулю, а вещественное устойчивое вспомогательное
поле даёт противоположный знак. Нужный симметричный средний модуль удалён
идеалом вырожденных двухформ.

\subsection{Фермионная мера}

Гауссов интеграл зависит от $\rho^2+r^2$ и не выводит мнимый контур
Хаббарда--Стратоновича. Для его переоткрытия потребовалось бы заранее
полученное четырёхфермионное ядро, которого в проекте нет.

Машинный реестр проверяет шесть одновременных требований для каждого
кандидата. Ни одна построенная ветвь не проходит их все.

\section{Отсутствует не один коэффициент, а типизированное соответствие}

В архитектуре $H_{15}$ первая калибровочно-инвариантная нейтринная масса
имеет высшую степень и схематически задаётся оператором Вайнберга
\begin{equation}
 \frac{c_{ij}}{\Lambda}
 (\overline{L_i^c}\,\widetilde H^*)
 (\widetilde H^\dagger L_j).
 \label{eq:v5-freeze-weinberg}
\end{equation}
Поэтому минимально отсутствующий объект нельзя описать словами
``ещё один скаляр''. Требуется майорановское соответствие высшей степени,
которое одновременно содержит:
\begin{enumerate}
 \item правильно представленное сечение или модуль с лептонным зарядом
 минус два;
 \item ковариантную производную с корневым или дефектным сектором;
 \item симметричную семейную карту спаривания, фиксированную до обращения к
 нейтринным данным;
 \item общий след или меру, нормирующие кинетический член, потенциал и
 фермионную связь;
 \item вариационно выведенную неоднородную конфигурацию и её длину
 локализации.
\end{enumerate}

Иными словами, открытая задача имеет тип
\begin{equation}
 \mathfrak P_\nu:
 \operatorname{Sym}^2(L_L\otimes H)
 \longrightarrow \mathcal E_{-2},
 \label{eq:v5-freeze-pairing-correspondence}
\end{equation}
а не тип свободно выбранного числа $m$ или функции $m(r)$.

\begin{theorem}[Заморозка части II]
Дефектно-транспортная ветвь Тома V построила содержательное кинематическое
и топологическое ядро, но не вывела локализованный нейтринный дефект.
Старая вершина $\Phi N^cN^c$ относится к иной архитектуре $H_{16}/M_{36}$.
Внутри сохранённого родителя $H_{15}/M_{35}$ минимальным объектом
переоткрытия является ещё не построенное майорановское соответствие высшей
степени \eqref{eq:v5-freeze-pairing-correspondence}.
\end{theorem}

\section{Архитектурная развилка}

После заморозки остаются два честных пути:
\begin{enumerate}
 \item сохранить $H_{15}$ и проверить существование соответствия
 \eqref{eq:v5-freeze-pairing-correspondence};
 \item перейти к $H_{16}$, зарегистрировать новый родитель $M_{36}$ и
 повторить гейты общего следа, меры, конденсата и нормировки.
\end{enumerate}
Правило минимального изменения выбирает первым путь $H_{15}$: его носитель
и веса уже выведены, тогда как второй путь является новой версией модели.

Следующий гейт должен классифицировать калибровочно- и
$C_3$-эквивариантные майорановские соответствия высшей степени на
$H_{15}$ и проверить, фиксируют ли существующие проектор и общий след их
семейный тензор без нового коэффициента.

\begin{protocol}
\begin{enumerate}
 \item кинематика переноса: \textbf{сохранена};
 \item голономный граничный нуль: \textbf{сохранён};
 \item локализованный объёмный дефект: \textbf{не выведен};
 \item совместимость старого $B-L$-вихря с $H_{15}$:
 \textbf{опровергнута};
 \item полный кандидат в существующем меню: \textbf{отсутствует};
 \item точный тип недостающего объекта: \textbf{зафиксирован};
 \item минимальная следующая архитектура: \textbf{высшая степень на $H_{15}$};
 \item физическое замыкание: \textbf{не достигнуто}.
\end{enumerate}
\end{protocol}

Машинный протокол сохранён в
\path{s2t_v5_defect_transport_status_freeze_gate_results.json}.